\subsection{C(d), continued: injectivity and the exact congruence fibres}

\begin{theorem}[C.13 --- Proved: real topological types remain distinct in the congruence image]
\label{C-real-type-injectivity}
Let $g\geq2$, and let $\tau$ and $\tau'$ be orientation-reversing
involutions of a closed oriented genus-$g$ surface. Conjugacy of their
outer actions on the profinite surface group implies conjugacy of their
mapping classes by an orientation-preserving mapping class.
Consequently the map from the real components of $\mathcal M_g$ to
conjugacy classes of involutions in the procongruence extended mapping
class group is injective. The same map to the full profinite extended
mapping class group is injective.

In genus three this distinguishes all six topological types: the
nondividing types with $r=0,1,2,3$ ovals and the dividing types with
$r=2,4$ ovals. The assertion concerns the images of these six types;
it does not classify the remaining profinite involutions.
\end{theorem}
\begin{proof}
Put $\Delta=\widehat{\pi_1(S_g)}$. This profinite surface group is
centre-free. We specify the extension recovered from an outer action
$\alpha:C_2\to\operatorname{Out}(\Delta)$:
\[
 E_\alpha=\operatorname{Aut}(\Delta)
       \mathbin{\times}_{\operatorname{Out}(\Delta)} C_2.
\]
For the arithmetic fundamental group $E$ of a real curve realizing
$\tau$, the map
\[
 E\longrightarrow E_\alpha,\qquad
 e\longmapsto(\operatorname{Ad}(e)|_\Delta,\overline e)
\]
is an isomorphism. Injectivity follows from the trivial centre of
$\Delta$. For surjectivity, choose any element of $E$ over the
specified element of $C_2$; its automorphism differs from the given
one by a unique inner automorphism. Multiplication by that element
of $\Delta$ supplies the preimage. This argument also applies to a
nonsplit real extension.

Choose an automorphism $b$ representing an outer conjugator between
$\alpha$ and $\alpha'$. It induces the extension isomorphism
\[
 E_\alpha\longrightarrow E_{\alpha'},\qquad
 (\varphi,\epsilon)\longmapsto
       (b\varphi b^{-1},\epsilon).
\]
Theorem~\ref{C-profinite-centralizer} therefore recovers the number
$r$ of real ovals as the number of geometric conjugacy classes of
sections. It recovers more than this number. For each section class
choose a representative $s_i$, and form
\[
 H_i=C_\Delta(s_i)\simeq\widehat{\mathbb Z},\qquad
 w_i=\operatorname{im}\bigl(
     1\in H_i/2H_i\longrightarrow H_1(\Delta,\mathbb F_2)\bigr).
\]
Here $1$ means the unique nonzero element of the group of order two,
so no choice of generator occurs. Changing the representative
conjugates the subgroup and leaves its image in homology unchanged.
Theorem~\ref{C-profinite-centralizer} identifies $w_i$ with the
modulo-two fundamental class of the corresponding oval $A_i$.
The extension isomorphism carries the unordered collection of these
classes to the corresponding collection for $\tau'$.

For $r>0$, the dividing invariant has the intrinsic expression
\[
 X(\mathbb C)\setminus X(\mathbb R)\text{ is disconnected}
 \quad\Longleftrightarrow\quad
 \sum_{i=1}^r w_i=0
       \quad\text{in }H_1(\Delta,\mathbb F_2).
\tag{C.13.1}
\]
Here is a direct topological proof. The complement of the oval union
$A$ is the orientation double cover of the interior of the connected
quotient surface, and thus has one or two components. With two
components, the closure of either component is a modulo-two
two-chain whose boundary is $A$. With one component, triangulate the
surface with $A$ as a subcomplex. A two-chain whose boundary is
supported on $A$ must have constant coefficient on all triangles
joined across edges outside $A$. Connectedness of the complement
forces one constant coefficient. Across every edge of $A$ the two
adjacent coefficients are then equal, so its boundary is zero.
It cannot have boundary $A$. This proves (C.13.1). When $r=0$ the
complement is connected and the type is already determined.

The extension isomorphism preserves both $r$ and (C.13.1).
The quotient-surface classification of
Proposition~\ref{C-classification}, or equivalently Gross--Harris,
Proposition 3.1 and its following realization statement
\cite{GrossHarris1981}, now gives an equivariant homeomorphism.
If it reverses orientation, compose it with the target involution
to obtain an orientation-preserving conjugator. Conjugacy in the
procongruence group is in particular conjugacy of these outer
actions. Conjugacy in the full completion maps to conjugacy in the
procongruence group, proving both injectivity assertions. Substituting
$g=3$ in the two quotient-surface formulas gives the six listed types.
\end{proof}


\begin{theorem}[C.14 --- Proved: the marked injectivity comparison, including relative boundary orientation]
\label{C-marked-injectivity}
For $g\geq2$, distinct real components of the moduli stack with
individually labelled real markings remain distinct in the
procongruence extended mapping class group and in its full profinite
completion. The analogous injectivity holds for unordered markings,
including conjugate pairs of nonreal markings.

More precisely, an isomorphism of the real arithmetic fundamental
group extensions which preserves the peripheral conjugacy classes
recovers the compact type $(g,r,a)$ and the real boundary incidence
data of Proposition~\ref{C-boundary-order}. In the dividing case,
the recovered cyclic orders on labelled marked ovals are identified
only by simultaneous reversal, as in
Proposition~\ref{C-marked-count}.
\end{theorem}
\begin{proof}
Write $U=X\setminus D$. Filling the punctures is the actual quotient
\[
 \pi_1(U_{\mathbb C})\longrightarrow\pi_1(X_{\mathbb C})
 =\pi_1(U_{\mathbb C})/
       \overline{\langle\!\langle I_d:d\in D\rangle\!\rangle}.
\tag{C.14.1}
\]
The corresponding quotient of the arithmetic extension recovers
the real extension of $X$. Hence the given isomorphism induces the
comparison in Theorem~\ref{C-real-type-injectivity}. In particular it
recovers $r$ and $a$. Proposition~\ref{C-boundary-order} recovers
the real punctures, their containing ovals, and the cyclic order
on each marked oval up to its separate reversal. Its graph is
understood as a graph with flags: the two flags of a cusp are the
two distinct elements of its local section torsor, even when they
are incident with the same vertex. Thus the one-edge loop and the
two-edge cycle retain their full incidence information.

We prove the additional orientation compatibility using the maps
of these graphs for finite covers. This supplies the specific
comparison between a discrete oval generator and its profinite
centralizer. Fix a marked oval $A$, an adjacent real interval,
and its section $\widetilde s$ of the extension for $U$. Let $s$
be its image in the extension for $X$, and put
\[
 P=\pi_1(X_{\mathbb C}),\qquad H=C_P(s)
       =\overline{\langle a\rangle}\simeq\widehat{\mathbb Z},
\]
where $a$ is one oriented circuit of the actual oval.
For an open normal subgroup $N$ of $P$ invariant under $s$, the
subgroup $N s(C_2)$ defines a finite real \'etale cover $X_N\to X$
with the chosen lift of the base real point. The section $s$ in
that subgroup selects the real oval $A_N$ through this lift.
The inverse image $U_N$ of $U$ has precisely the preimages of $D$
as punctures, and the section $\widetilde s$ selects the base
interval of the incidence cycle on $A_N$.

We record the fibre and monodromy identifications explicitly.
With the chosen sheet denoted by $1$, the complex fibre of this
cover is $P/N$. Its real points are the fixed cosets for $s$.
The fibre on the selected connected real oval is the monodromy
orbit of $1$ under the discrete oval loop:
\[
 \langle a\rangle N/N=HN/N\simeq H/(H\cap N).
\tag{C.14.2}
\]
The first equality follows by lifting paths along the real circle;
the second follows by taking the closure in this finite fibre.
In particular the real oval cover has degree
$[H:H\cap N]$. For $h\in H$, normality of $N$ and the equality
$hs=sh$ show that $h$ normalizes $Ns(C_2)$, giving an actual
real deck transformation. Its action on this selected component
agrees with (C.14.2). Equivalently,
its inverse limit is the centralizer: the inverse limit of all
fixed-coset sets is $P^s=H$, and (C.14.2) identifies the chosen
monodromy orbits with its finite cyclic quotients. These quotients
are cofinal because Theorem~\ref{C-profinite-centralizer} proves
the full profinite topology on $\langle a\rangle$.

The graph construction is functorial for these covers, not merely
for their individual graphs. A cusp maps to its image cusp; each
of its two local section classes maps to the corresponding local
class; a global section class maps by composition with the covering
homomorphism. These maps commute with the local-to-global incidence
maps. Since $X_N\to X$ is \'etale also at the filled punctures,
the result on the selected cycles is the degree-$[H:H\cap N]$
cover of the base cycle, with its flags and its selected sheet.
The deck actions in (C.14.2) act on these graph covers. The quotient
map (C.14.1) ensures that a peripheral extension isomorphism sends
exactly this system of compact \'etale covers, and their inverse
images in $U$, to the corresponding system on the other curve.
It therefore gives compatible isomorphisms of the graph-cover
maps and their deck actions.

Choose orientations of the source and target base cycles. The
base graph isomorphism has degree $\epsilon\in\{1,-1\}$. Lifting
one positively oriented circuit in every graph cover now gives
\[
 b(a)\equiv (a')^{\epsilon}\pmod {H'\cap b(N)}.
\]
The same sign occurs at every level because each graph-cover
isomorphism lies over the same base graph isomorphism. Cofinality
gives the equality $b(a)=(a')^{\epsilon}$ in $H'$, after the
conjugations used to identify the chosen sections. An arbitrary
unit of $\widehat{\mathbb Z}$ is thus excluded by the actual
graph-cover maps. On homology the equality is independent of
those conjugations.

Suppose now that $X$ is dividing. Orient its $r$ ovals as the
boundary of one half of $X(\mathbb C)$. Their classes satisfy
exactly the relation
\[
 \ker\left(\bigoplus_{i=1}^r\widehat{\mathbb Z}[a_i]
      \longrightarrow H_1(X(\mathbb C),\widehat{\mathbb Z})\right)
     =\widehat{\mathbb Z}(1,\ldots,1).
\tag{C.14.3}
\]
Indeed, in the integral relative homology sequence of the pair
$(X(\mathbb C),\bigcup A_i)$, the relative two-dimensional group
is generated by the two halves. Their boundary classes are
$\sum_i[a_i]$ and its negative. Exactness gives the stated
kernel over $\mathbb Z$, and tensoring with the flat
$\mathbb Z$-module $\widehat{\mathbb Z}$ gives (C.14.3).

An isomorphism sends a generator of an unmarked oval centralizer
to a unit power of a generator of the corresponding oval. On a
marked oval the preceding graph argument identifies that unit
with its graph sign $\epsilon_i\in\{1,-1\}$. Applying the
homology isomorphism to (C.14.3) shows that all these unit
multipliers are equal: its image vector must lie in the target
diagonal kernel. In particular all marked-oval signs agree.
The statement is also valid when there is only one marked oval;
then no relative sign needs to be compared. This proves exactly
the simultaneous-reversal equivalence for labelled dividing
components. For nondividing components the independent reversals
are already the equivalence in
Proposition~\ref{C-marked-count}.

For unordered markings, the extension also gives the two-element
orbits of nonreal peripheral classes and hence the number of
interior marked points of the quotient surface. The real incidence
cycles give the numbers of boundary markings on each boundary
component. Configurations of a fixed number of unordered interior
points in a connected surface are connected, as are the boundary
configurations with fixed cyclic incidence. Homeomorphisms of the
quotient surface realize the boundary permutations, and its double
gives the equivariant homeomorphism upstairs. These data therefore
give exactly the unordered topological type. The classification
in Proposition~\ref{C-marked-count} and the full-completion to
procongruence comparison finish the two injectivity assertions.
\end{proof}

\begin{proposition}[C.15 --- Proved: the congruence-fibre action and centralizer image]
\label{C-congruence-fibre-action}
Write
\[
 1\longrightarrow K\longrightarrow P^{\pm}
       \xrightarrow{q} Q^{\pm}\longrightarrow1
\]
for the full-completion to procongruence quotient of an extended
mapping class group, and let $P$ and $Q$ denote its
orientation-preserving subgroups. Fix a discrete real involution
$\tau\in P^{\pm}$, put $\bar\tau=q(\tau)$, and let
$\alpha=\operatorname{Ad}(\tau)$. The action
\[
 C_Q(\bar\tau)\times H^1(C_2,K)\longrightarrow H^1(C_2,K),
 \qquad q(g)\star[k]=[gk\alpha(g)^{-1}]
\tag{C.15.1}
\]
is well defined, where $g\in P$ lifts the indicated centralizer
element. The fibre over $[\bar\tau]$ of the map on conjugacy
classes of orientation-reversing involutions is exactly
\[
 C_Q(\bar\tau)\backslash H^1(C_2,K).
\tag{C.15.2}
\]
For every cocycle $k\in K$, the centralizer fits into the exact
sequence
\[
 1\longrightarrow K^{k\tau}\longrightarrow C_P(k\tau)
 \xrightarrow{q}
 \operatorname{Stab}_{C_Q(\bar\tau)}([k])
 \longrightarrow1.
\tag{C.15.3}
\]
In particular the image on centralizers at $\tau$ is the
stabilizer of the trivial cohomology class for (C.15.1).
For the general extension calculation underlying these formulas,
the action does not automatically preserve that distinguished
class, as the finite example in the proof verifies.
\end{proposition}
\begin{proof}
The elements of order two above the particular element
$\bar\tau$ are $k\tau$ with $k\alpha(k)=1$.
Conjugation by $b\in K$ changes $k$ to
$bk\alpha(b)^{-1}$, giving precisely the stated
nonabelian $H^1$. If $q(g)$ centralizes $\bar\tau$, then
$gk\alpha(g)^{-1}$ lies in $K$, and
\[
 g(k\tau)g^{-1}=gk\alpha(g)^{-1}\tau.
\]
This proves that it is again a cocycle. Changing the lift to
$bg$ changes the result by the displayed $K$-conjugacy relation.
Changing a representative of $[k]$ has the same effect, and
composition of conjugations proves the action law.

An involution whose image is $Q$-conjugate to $\bar\tau$ may
first be conjugated by a lift in $P$ so that its image equals
$\bar\tau$. Two such representatives are $P$-conjugate
precisely when their $K$-conjugacy classes are in the same orbit
of (C.15.1). This proves (C.15.2).

The kernel of the centralizer map is exactly $K^{k\tau}$.
Its image fixes $[k]$. Conversely, take an element of $C_Q(\bar\tau)$
fixing $[k]$, and choose a lift $g\in P$. By the definition of
the cohomology class, there is $b\in K$ with
$g(k\tau)g^{-1}=b(k\tau)b^{-1}$. Then $b^{-1}g$ centralizes
$k\tau$ and has the prescribed image in $Q$. This proves
surjectivity onto the stated stabilizer and all of (C.15.3).

To verify the last qualification, take the finite extension with
$P^{\pm}=D_8=\langle r,t\mid r^4=t^2=1,\ trt=r^{-1}\rangle$,
$P=\langle r\rangle$, $K=\langle r^2\rangle$, and $\tau=t$.
Here $Q=P/K$ has order two, and the whole $Q$ centralizes
$\bar\tau$. The action of $\tau$ on $K$ is trivial, so
$H^1(C_2,K)=\{[1],[r^2]\}$. The element $q(r)$ sends $[1]$
to $[r\alpha(r)^{-1}]=[r^2]$. Thus (C.15.2) has one orbit,
whereas this $H^1$ has two elements, and
$C_P(t)=K$ has trivial image in $Q$. These explicit groups
verify why the orbit and stabilizer in the formulas must be
retained. They assert no nonvanishing congruence kernel for a
mapping class group.
\end{proof}

\begin{remark}[Scope of the C.13--C.15 advance]
Theorems~\ref{C-real-type-injectivity} and
\ref{C-marked-injectivity} establish injectivity for every real
topological type under discussion. Formula (C.15.2) records the
entire remaining fibre over each known congruence class, and
(C.15.3) records its exact centralizer map. The cardinalities of
these orbit sets, the possible additional involution classes in
$Q^{\pm}$, and the comparison of $C_P(\tau)$ with the closure
of the discrete mapping-class centralizer remain uncomputed in
genus three and higher. Thus the surjectivity and complete
centralizer request in Proposition~\ref{C-high-genus-open}
remains incomplete. No assertion of mapping-class-group goodness
or of a vanishing congruence kernel is used in this advance.
\end{remark}

\begin{proposition}[C.16 --- Proved: the hyperelliptic part of every genus-three real centralizer]
\label{C-hyperelliptic-reflections}
Fix a hyperelliptic involution $\iota$ of $S_3$, and put
$H^{\pm}=C_{\Gamma_3^{\pm}}(\iota)$ and
$H=H^{\pm}\cap\Gamma_3$. The map of full profinite completions
\[
 \widehat{H^{\pm}}\longrightarrow\widehat{\Gamma_3^{\pm}}
\tag{C.16.1}
\]
is injective, with image $\overline{H^{\pm}}$. Every
orientation-reversing involution of this closed subgroup is
$\overline H$-conjugate to a discrete involution in $H^{\pm}$.
For every such discrete involution $\tau$,
\[
 C_{\widehat{\Gamma_3}}(\tau)\cap\overline H
       =C_{\overline H}(\tau)
       =\overline{C_H(\tau)}.
\tag{C.16.2}
\]
Each of the six real genus-three topological types has a
representative $\tau$ in this $H^{\pm}$.
\end{proposition}
\begin{proof}
The hyperelliptic congruence subgroup theorem, in the precise
form of Boggi, Theorem 7.1~\cite{CextraBoggiHyperelliptic}, says
that the composite
\[
 \widehat H\longrightarrow\widehat{\Gamma_3}
       \longrightarrow\operatorname{Out}
               (\widehat{\pi_1(S_3)})
\]
is injective. Its first map is therefore injective. The
orientation subgroups have index two, so the kernel of
(C.16.1) lies in $\widehat H$ and is trivial. Compactness and
density identify the image with the indicated closure.

We give the goodness argument for the group actually used here.
The branched double cover supplies the Birman--Hilden sequence
\[
 1\longrightarrow\langle\iota\rangle\longrightarrow H
       \longrightarrow\Gamma_{0,8}\longrightarrow1.
\tag{C.16.3}
\]
The map sends a commuting mapping class to the induced mapping
class of the quotient sphere with its eight branch points.
Every such sphere mapping class lifts: the double cover on the
punctured sphere is the one sending every peripheral meridian
to the nonidentity element of $C_2$, and that homomorphism is
preserved by every permutation of the branch points. The kernel
in (C.16.3) is the deck group; the mapping-class exactness is the
Birman--Hilden sequence (3) of the cited primary source.

Intersect the inverse image of the pure sphere group with the
kernel of $H\to\operatorname{Sp}_6(\mathbb F_3)$. This is a
finite-index subgroup $L$ of $H$ which misses $\iota$, since
$\iota$ acts by $-1$ on first homology. By (C.16.3), $L$ is
isomorphic to a finite-index subgroup of the pure group
$\Gamma_{0,8}^{\mathrm{pure}}$. The latter is good, as can be
proved directly from its configuration-space fibrations.
Fix three points at $0,1,\infty$; forgetting the last of the
remaining ordered points has a punctured-plane fibre and a
continuous section: choose the additional real point
$2+\sum_j|z_j|$. Thus the fundamental-group extension is a
split extension by a finitely generated free group. Every
finite-index subgroup of this free kernel contains a
characteristic finite-index subgroup. The action on the
corresponding finite quotient has finite image, giving a
finite quotient of the whole semidirect product which
restricts to that quotient of the kernel. Its induced
profinite topology is therefore the full one, and the
completed extension is exact.

Free groups are good; induction on the number of points,
using the two Hochschild--Serre spectral sequences, proves
goodness of the pure sphere group. The finite-type free
resolutions give finite cohomology with finite coefficients
and a finite cohomological dimension at every stage.
Goodness passes to finite-index subgroups by Shapiro's lemma,
and to finite extensions by taking a normal finite-index
subgroup and comparing the Hochschild--Serre sequences.
It follows successively for $L$, $H$, and $H^{\pm}$.
This also proves finite virtual two-cohomological type.
These groups have finite-type resolutions in each degree,
so passage from finite modules to arbitrary discrete
continuous $\mathbb F_2$-modules is obtained by filtered
unions, as in Theorem~\ref{C-profinite-centralizer}.

Residual finiteness follows from the faithful outer action
just used, with orientation detected on the fundamental
class. All hypotheses of Boggi--Zalesskii, Theorem A, now
hold for $H^{\pm}$ and the prime two~\cite{BoggiZalesskii2024}.
That theorem gives the asserted conjugacy and centralizer
identities inside $\widehat{H^{\pm}}$. An
orientation-reversing conjugator can be made
orientation-preserving by multiplying it by the source
involution. Intersecting the centralizer identity with the
open subgroup $\widehat H$ gives (C.16.2).

For completeness, the following smooth hyperelliptic models
give all six real types; $a=0$ denotes dividing and $a=1$
nondividing:
\[
\begin{array}{c|c|c}
 y^2=f(x)&r&a\\ \hline
 -(x^8+1)&0&1\\
 (x^2-1)(x^6+1)&1&1\\
 (x^2-1)(x^2-4)(x^4+1)&2&1\\
 (x^2-1)(x^2-4)(x^2-9)(x^2+1)&3&1\\
 \prod_{j=1}^{4}(x^2-j^2)&4&0\\
 x^8+1&2&0
\end{array}
\tag{C.16.4}
\]
Every polynomial is square-free of degree eight, so its smooth
completion has genus three. With $2k>0$ real branch points
and positive leading coefficient, the $k$ positive intervals
on the real projective line lift to $k$ ovals. For
$k=1,2,3$, the preimage of each open half-plane is connected
because it contains a branch point. A negative real interval
joins these two preimages through nonreal points of the curve;
hence the complement of the ovals is connected. The case
$k=4$ has four ovals and is dividing by
Proposition~\ref{C-classification}.

The negative polynomial in the first row has no real points,
including at infinity. For $x^8+1$, the real double cover of
$\mathbb P^1(\mathbb R)$ is the union of two circles: in the
infinity chart $u=1/x$, $v=y/x^4$, each sign of $y$ joins to
the same sign of $v$. Every point above the real projective
line is real. The preimages of the upper and lower half-planes
are each connected and form the two components of its
complement, proving the last row. Complex conjugation in
each model commutes with $(x,y)\mapsto(x,-y)$. All genus-three
hyperelliptic involutions are conjugate, by the branched
double-cover construction with eight branch points. Transport
to the fixed $\iota$ therefore puts each of the six real
involutions in $H^{\pm}$.

The computed group in (C.16.2) is the specified intersection
with $\overline H$. This calculation does not establish that
every ambient profinite reflection belongs to a conjugate
of $\overline{H^{\pm}}$, or that its entire centralizer lies
in this subgroup. Those are additional, unresolved steps
in the full request.
\end{proof}

\begin{theorem}[C.17 --- Proved: congruence centralizers for the five genus-three types with real points]
\label{C-genus-three-centralizer-reduction}
Let $\tau\in\Gamma_3^{\pm}$ be a real involution with at least
one oval, and put
\[
 P=\widehat{\Gamma_3},\qquad Q=\check\Gamma_3,
 \qquad K=\ker(P\xrightarrow{q}Q),\qquad
 C=C_{\Gamma_3}(\tau).
\]
For closures in the indicated groups, restriction of $q$ gives
an isomorphism
\[
 \overline C^{\,P}\xrightarrow{\ \sim\ }
 C_Q(\bar\tau)=\overline C^{\,Q}.
\tag{C.17.1}
\]
The full profinite centralizer has the topological semidirect
product description
\[
 C_P(\tau)=K^\tau\rtimes\overline C^{\,P}.
\tag{C.17.2}
\]
Thus the centralizer action (C.15.1) fixes the trivial class
for each of these five genus-three types. The value of the
remaining group $K^\tau$ and the other cohomology orbits is
not determined here. The fixed-point-free genus-three type
is excluded from this theorem.
\end{theorem}
\begin{proof}
Let $A=\{A_1,\ldots,A_r\}$ be the real oval multicurve and
$\Delta=\widehat{\pi_1(S_3)}$. The ovals are essential and
pairwise nonisotopic. Indeed a contractible oval would bound
a disc exchanged with the other side, forcing a sphere;
two parallel ovals would give an annular component of the
complement, forcing an annulus on each side in the dividing
case or the whole complement to be that annulus in the
nondividing case. Both alternatives force genus one.

First we prove the precise preservation needed to apply a
multicurve stabilizer theorem. Take $B\in C_Q(\bar\tau)$.
The pullback description in
Theorem~\ref{C-real-type-injectivity} shows that $B$ permutes
the section classes of the real surface extension and hence
their centralizers $\overline{\langle a_i\rangle}$. To identify
the actual curve classes rather than just their procyclic
subgroups, choose an auxiliary real marked point on $A_i$.
The procongruence Birman comparison is the isomorphism
\[
 \check\Gamma_{3,1}^{\pm}
 \xrightarrow{\ \sim\ }
 \operatorname{Aut}(\Delta)
       \mathbin{\times}_{\operatorname{Out}(\Delta)}
                         \check\Gamma_3^{\pm}.
\tag{C.17.3}
\]
Here the map on the left is conjugation on the point-pushing
subgroup $\Delta$. Its faithfulness and its exact kernel and
quotient are respectively Corollaries 4.5 and 4.7 of
Boggi~\cite{CextraBoggiCongruence2020}, for the class of all
finite groups. The kernel maps isomorphically to
$\operatorname{Inn}(\Delta)$ because $\Delta$ is centre-free.
Surjectivity onto the displayed pullback follows by lifting
the outer element and then multiplying by the unique needed
inner automorphism. Adjoining the actual pointed reflection
gives the same isomorphism for the orientation extensions.

Choose a lift of $B$ in (C.17.3). Its conjugate of the pointed
reflection is a section of the surface extension and therefore,
by Theorem~\ref{C-profinite-centralizer}, is $\Delta$-conjugate
to a pointed reflection on one of the actual ovals, say $A_j$.
Adjust the lift by that element of $\Delta$. Injectivity in
(C.17.3) shows that it conjugates the two pointed reflections
as elements of the punctured procongruence group, not merely
their filled automorphisms. Apply the one-marked-oval graph
argument of Theorem~\ref{C-marked-injectivity} to its action
on the punctured surface group. It gives
\[
 B([a_i])=[a_j^{\epsilon_i}],\qquad
                   \epsilon_i\in\{1,-1\},
\tag{C.17.4}
\]
where brackets mean geometric conjugacy classes. Repeating
this construction for each oval shows that $B$ preserves
the actual unoriented multicurve $A$ in the procongruence
curve complex.

Let $L=\operatorname{Stab}_{\Gamma_3}(A)$ and let $L_0$ be
the finite-index subgroup fixing each oval and its orientation.
Let $S_1,\ldots,S_d$ be the components after cutting along
$A$ and replacing their boundary circles by labelled
punctures. The ordinary cutting homomorphism has exact
sequence
\[
 1\longrightarrow T\longrightarrow L_0
       \longrightarrow\prod_{j=1}^{d}\Gamma(S_j)^{\mathrm{pure}}
       \longrightarrow1,\qquad T\simeq\mathbb Z^r,
\tag{C.17.5}
\]
where $T$ is generated by the individual oval Dehn twists.
Theorem 4.5 of Boggi~\cite{CextraBoggiProcongruence2014}
identifies $Q_A$ with $\overline L^{\,Q}$ and gives the exact
completed cutting sequence
\[
 1\longrightarrow\widehat T\longrightarrow Q_{\vec A}
       \longrightarrow\prod_{j=1}^{d}
                  \check\Gamma(S_j)^{\mathrm{pure}}
       \longrightarrow1.
\tag{C.17.6}
\]
In particular the twist subgroup here is the injected full
$\widehat{\mathbb Z}^{r}$; it is not a quotient of that group.

Every component $S_j$ in this genus-three situation has genus
at most two. The complete list is
\[
\begin{array}{c|c|c}
 (r,a)&S_3\setminus A\text{ with boundaries punctured}
       &\text{component genera}\\ \hline
 (1,1)&S_{2,2}&2\\
 (2,1)&S_{1,4}&1\\
 (3,1)&S_{0,6}&0\\
 (2,0)&S_{1,2}\amalg S_{1,2}&1,1\\
 (4,0)&S_{0,4}\amalg S_{0,4}&0,0
\end{array}
\tag{C.17.7}
\]
The genus-at-most-two congruence subgroup theorem therefore
identifies each right-hand group in (C.17.6) with the full
completion of its group in (C.17.5). For the positive-genus
component groups this is the precise range of Theorem 1.4(ii)
of \cite{CextraBoggiCongruence2020}; the genus-zero component
groups are included in Theorem 7.1 of
\cite{CextraBoggiHyperelliptic} with the underlying closed
surface a sphere. For the class of all finite groups, the
level-two isomorphism in the first source extends to the
whole group by its finite level quotient.

Completing (C.17.5) is exact at its middle and right terms.
The map of its completed kernel to (C.17.6) is the identity
on $\widehat T$, so it is also injective on the left. The
resulting diagram of exact sequences has isomorphisms at
both ends; its middle map is an isomorphism. Thus
$\widehat{L_0}\simeq Q_{\vec A}$. The finite permutation and
orientation action on the ovals has the same image for
$L$ and its closure, which gives
\[
 \widehat L\simeq Q_A.
\tag{C.17.8}
\]
It follows also that $\widehat L\to P$ is injective, since
its composite with $q$ is the injection in (C.17.8).

We check goodness for the stabilizer to which we will apply
the finite-subgroup theorem. Pure mapping class groups of
genus zero are good by the configuration argument in
Proposition~\ref{C-hyperelliptic-reflections}. The group
$\Gamma_{1,1}\simeq\operatorname{SL}_2(\mathbb Z)$ is
virtually free. The group $\Gamma_2$ is hyperelliptic, so
the same finite-index sphere argument proves its goodness.
Adding punctures preserves goodness by the exact completed
Birman sequence and the goodness of surface groups proved
in Theorem~\ref{C-profinite-centralizer}, using the
Hochschild--Serre comparison. These arguments also give
finite virtual cohomological type and finite-type resolutions.
The product of these component groups has the same
properties. The exactness just verified for the completion
of (C.17.5), and goodness of $T\simeq\mathbb Z^r$, now give
goodness of $L_0$ by the same spectral-sequence comparison.
Finite extension gives goodness and finite virtual
two-cohomological type for $L$ and for
$L^{\pm}=L\rtimes\langle\tau\rangle$. The groups are
residually finite through their faithful congruence actions.
The finite-type resolutions justify the discrete-module
form of the cohomology comparison required in Theorem A.

Because $\tau$ preserves $A$, (C.17.8) extends to
$\widehat{L^{\pm}}\simeq Q_A\rtimes\langle\bar\tau\rangle$.
Boggi--Zalesskii, Theorem A~\cite{BoggiZalesskii2024}, applied
to this particular good group and the prime two, yields
\[
 C_{Q_A}(\bar\tau)=\overline{C_L(\tau)}^{\,Q}.
\]
By (C.17.4), the entire centralizer in $Q$ lies in $Q_A$.
The same preservation statement holds for the discrete
centralizer, so $C_L(\tau)=C$. This proves the centralizer
identity in (C.17.1). The injection $\widehat L\to P$ and
its isomorphism onto $Q_A$ identify the two closures of
$C$, proving all of (C.17.1).

Finally, $\overline C^{\,P}$ centralizes $\tau$, and maps
isomorphically onto the full quotient centralizer. For an
element $u\in C_P(\tau)$, multiply $u$ by the inverse of
the unique element of $\overline C^{\,P}$ with image $q(u)$.
The result lies in $K$ and still centralizes $\tau$.
This factorization is unique because the two factors
intersect trivially. Conjugation by the second factor
preserves $K^\tau$, and compactness gives the asserted
topological semidirect product. The final assertion about
the trivial class follows from (C.15.3) and the now
surjective centralizer map.
\end{proof}

\begin{proposition}[C.18 --- Proved: the fixed-point-free comparison]\label{C-free-reflection-comparison}
Let $\tau$ be a fixed-point-free orientation-reversing
involution of $S_3$, let $N_4=S_3/\langle\tau\rangle$, and put
\[
 E=\widehat{\pi_1(N_4)},\qquad
 \Delta=\widehat{\pi_1(S_3)},\qquad
 Q^\pm=\check\Gamma_3^\pm,\qquad Q=\check\Gamma_3.
\]
The inclusion of the orientation double cover gives the
exact sequence $1\to\Delta\to E\to C_2\to1$, whose outer
action is $\bar\tau$.
Let $\operatorname{Aut}_Q(E;\Delta)$ be the group of
automorphisms preserving $\Delta$ whose restriction to
$\Delta$ has outer class in $Q^\pm$, and define
\[
 \mathcal A=\operatorname{Aut}_Q(E;\Delta)/
                         \operatorname{Inn}(E).
\]
There is a canonical topological isomorphism
\[
 \mathcal A
 \ \xrightarrow{\ \sim\ }\
 C_{Q^\pm}(\bar\tau)/\langle\bar\tau\rangle
 \ \xrightarrow{\ \sim\ }\
 C_Q(\bar\tau).
\tag{C.18.1}
\]
The first arrow is induced by restriction. The second
sends the class of $b$ to $b$ when $b$ preserves orientation
and to $b\bar\tau$ otherwise.
Under this isomorphism the closure of the ordinary mapping
class group of $N_4$ is exactly
$\overline{C_{\Gamma_3}(\tau)}^{\,Q}$.
Consequently the remaining congruence-centralizer assertion
for this type is the explicit density problem
\[
 \overline{\operatorname{Mod}(N_4)}^{\,\mathcal A}
                       \stackrel{?}{=}\mathcal A.
\tag{C.18.2}
\]
This density, and the corresponding full-completion
centralizer, are not established here.
\end{proposition}
\begin{proof}
The index-two subgroup $\pi_1(S_3)$ has its full profinite
topology in $\pi_1(N_4)$, so completion gives the asserted
extension. Its action on the completed orientation class
of $\Delta$ is negative; hence $\bar\tau$ is a nontrivial
outer involution.

Choose $e\in E\setminus\Delta$ and write
$a=\operatorname{Ad}(e)|_\Delta$ and $z=e^2\in\Delta$.
Thus $a^2=\operatorname{Inn}(z)$.
For $b\in\operatorname{Aut}(\Delta)$ whose outer class
centralizes $\bar\tau$, there is a unique $d\in\Delta$
such that
\[
 b a b^{-1}=\operatorname{Inn}(d)a.
\]
Uniqueness uses $Z(\Delta)=1$. Squaring this identity
gives
$\operatorname{Inn}(b(z))
 =\operatorname{Inn}(d\,a(d)\,z)$, whence
\[
 b(z)=d\,a(d)\,z.
\]
These two identities show that $b$ extends to an
automorphism of $E$ by the exact formula $e\mapsto de$:
they preserve, respectively, conjugation on $\Delta$
and the square of $e$. Since $E$ is the disjoint union
of $\Delta$ and $\Delta e$, the formula is continuous
and bijective on these two cosets. The extension is
unique, since an automorphism restricting to the identity
on $\Delta$ would have $e\mapsto de$ with
$d\in Z(\Delta)$.
This proves that restriction identifies
$\operatorname{Aut}_Q(E;\Delta)$ with the inverse image
of $C_{Q^\pm}(\bar\tau)$ in
$\operatorname{Aut}(\Delta)$.

The subgroup of inner automorphisms induced by $\Delta$
restricts to $\operatorname{Inn}(\Delta)$; conjugation
by $e$ contributes precisely the outer class $\bar\tau$.
Quotienting therefore gives the first isomorphism in
(C.18.1). The orientation homomorphism and the central
element $\bar\tau$ give the direct product
\[
 C_{Q^\pm}(\bar\tau)
       =C_Q(\bar\tau)\times\langle\bar\tau\rangle,
\]
which proves the displayed formula for the second arrow.
All groups here are closed subgroups or quotients of
automorphism groups of finitely generated profinite
groups. The continuous bijections are thus
homeomorphisms.

For the discrete comparison, every homeomorphism of
$N_4$ has two lifts to its orientation double cover and
exactly one preserves orientation. This choice is
compatible with composition and isotopy. The induced
map
$\operatorname{Mod}(N_4)\to C_{\Gamma_3}(\tau)$ is
an isomorphism: the quotient map from the extended
centralizer has kernel $\langle\tau\rangle$ and is onto.
This is the precise orientation-double-cover statement
recorded by Stukow~\cite[p.~1030]{CextraStukow2016}.
The induced action on $\pi_1(N_4)$, followed by completion
and (C.18.1), is exactly the action of that positive
lift on $\Delta$. Taking closures proves the final
assertion and identifies the unproved equality
(C.18.2) without asserting it.
\end{proof}

\begin{proposition}[C.19 --- Proved: no finite-index hyperelliptic reduction for the free genus-three type]
\label{C-no-finite-index-hyperelliptic}
Let $\tau$ be the fixed-point-free orientation-reversing involution
of $S_3$, and let
\[
 j:\operatorname{Mod}(N_4)
      \xrightarrow{\;\sim\;}C_{\Gamma_3}(\tau)
\]
be the unique orientation-preserving lift of the orientation double
cover, as in Proposition~\ref{C-free-reflection-comparison}.
For every finite-index subgroup $H$ of $\operatorname{Mod}(N_4)$,
the centralizer $C_{\Gamma_3}(j(H))$ has no nonidentity element of
finite order. In particular no such $j(H)$ is contained in the
centralizer of a hyperelliptic involution of $S_3$.
The standard hyperelliptic mapping class subgroup of
$\operatorname{Mod}(N_4)$ therefore has infinite index.
\end{proposition}
\begin{proof}
We give the surface and all comparison maps explicitly. Glue two
pairs of pants $P,Q$ along three pairs of boundary circles
$a,b,c$. Choose the gluing signs so that the orientation double
has the following dual graph:
\[
\begin{array}{c|c}
\text{lifted curve}&\text{incident pairs of pants}\\ \hline
a_+,b_+&P_+,Q_+\\
a_-,b_-&P_-,Q_-\\
c_+&P_+,Q_-\\
c_-&P_-,Q_+
\end{array}
\tag{C.19.1}
\]
Equivalently the voltage of the first two edges of the base
theta graph is zero and that of the third is one in $C_2$.
The cover is connected: the cycle using $a$ and $c$ has nonzero
voltage. Thus the base is nonorientable. It is closed and has
Euler characteristic $-2$, so it is $N_4$; the connected
orientation cover has Euler characteristic $-4$ and is $S_3$.
Each of $a,b,c$ is two-sided, with an annular neighbourhood
arising from the gluing collars. The six lifted curves form a
full pants decomposition of $S_3$. They are essential and
pairwise nonisotopic: a disk or an annulus bounded by some
of them would be a union of complementary pants, contradicting
its Euler characteristic.

Choose an orientation on each base annulus and let
$t_a,t_b,t_c$ be the resulting twists. The orientation-cover map
restricts to a homeomorphism on each lifted annulus. Its deck
involution reverses the surface orientation and exchanges the
two lifts. Consequently the positive-lift homomorphism gives
the exact identities
\[
 j(t_a)=T_{a_+}T_{a_-}^{-1},\qquad
 j(t_b)=T_{b_+}T_{b_-}^{-1},\qquad
 j(t_c)=T_{c_+}T_{c_-}^{-1},
\tag{C.19.2}
\]
after choosing the indicated labels using those annular
orientations. The six Dehn twists generate the injected group
$\mathbb Z^6$. The three displayed elements generate its
subgroup with exponent vectors
$(u,-u,v,-v,w,-w)$, hence an injected $\mathbb Z^3$.
Here we use the uniqueness of the support and exponents of a
multitwist, including arbitrary nonzero signed exponents;
the precise statement is Farb--Margalit, Lemma 3.17
\cite{CextraFarbMargalit}.

Finite index supplies positive integers $m_a,m_b,m_c$ such that
$t_a^{m_a},t_b^{m_b},t_c^{m_c}\in H$.
Let $\iota\in C_{\Gamma_3}(j(H))$ have finite order.
Conjugation by an orientation-preserving mapping class takes
$T_d$ to $T_{\iota(d)}$ with the same exponent.
Applying multitwist uniqueness to each identity in (C.19.2)
therefore shows that $\iota$ fixes each of
$a_+,a_-,b_+,b_-,c_+,c_-$ individually: the two exponents in
one pair are $m$ and $-m$, so they cannot be interchanged.

Every vertex in (C.19.1) has a different set of incident labelled
edges. An automorphism fixing all six edges consequently fixes
all four vertices. A representative of $\iota$ preserving the
multicurve thus preserves each complementary pair of pants
and each of its three boundary circles. Its restriction, with
boundaries replaced by punctures, lies in
$\operatorname{PMod}(S_{0,3})=1$.
The cutting homomorphism has exactly the multitwist kernel
$\mathbb Z^6$: these are Proposition 2.3 and Proposition 3.20
of the same primary text, with its explicit cutting map
\cite{CextraFarbMargalit}. Thus $\iota$ belongs to this
torsion-free kernel and is the identity.

This excludes every nonidentity finite-order positive class,
without an assumption that it already commutes with $\tau$.
For the standard hyperelliptic subgroup of $\operatorname{Mod}(N_4)$,
the positive-lift homomorphism is explicitly a map into the
genus-three hyperelliptic centralizer
\cite[p.~1030]{CextraStukow2016}. It cannot have finite-index
domain by the result just proved. The finite graph and the
three independent exponent vectors are also checked by
\texttt{code/B/n4\_pants\_check.py}; the infinite-group
conclusion follows from the written cutting argument.
\end{proof}
